% !TeX program = xelatex
\documentclass[11pt,a4paper]{ctexart}

\usepackage[left=2.4cm,right=2.4cm,top=2.1cm,bottom=2.1cm]{geometry}
\usepackage{amsmath,amssymb,mathtools,braket}
\usepackage{xcolor}
% inline 选项用于启用 enumerate* 行内列表环境。
\usepackage[inline]{enumitem}
\usepackage{fancyhdr}

% ============================================================
% 基本信息：只需修改下面几项
% ============================================================
\newcommand{\coursename}{高等量子力学}
\newcommand{\assignmentnumber}{第 1 次作业}
\newcommand{\teachername}{\underline{\hspace{2.8cm}}}
\newcommand{\studentname}{\underline{\hspace{2.8cm}}}
\newcommand{\studentid}{\underline{\hspace{3.0cm}}}

% ============================================================
% 数学公式书写规范
% ============================================================
% 1. 行内公式使用一对美元符号。
% 2. 独立成行的公式统一使用 equation 环境，不使用反斜杠方括号形式。
% 3. 虚数单位 i、自然常数 e 和微分符号 d 均使用直立体 \mathrm。
%

% 独立公式的推荐写法：
% \begin{equation}
%   ...
% \end{equation}
%

% ============================================================
% 页面样式
% ============================================================
\setlength{\parindent}{2em}
\setlength{\parskip}{0.25em}
\setlength{\headheight}{14pt}
\pagestyle{fancy}
\fancyhf{}
\lhead{\coursename}
\rhead{\assignmentnumber}
\cfoot{\thepage}
\renewcommand{\headrulewidth}{0.4pt}

% ============================================================
% 列表参数
% ============================================================
% enumerate 是每个条目 另起一行 的普通编号列表。
% label 决定编号格式：\alph* 依次产生 a、b、c，外加括号后显示为 (a)、(b)、(c)。
\setlist[enumerate,1]{
  label=(\alph*),
  leftmargin=3.2em,
  itemsep=0.55em,
  topsep=0.45em
}

% inlineenum 是基于 enumerate* 定义的行内编号列表。
% 所有格式均在此处统一设置，正文中不必为每一道题重复填写参数。
% label 决定编号格式；labelsep 决定编号与本项正文之间的距离。
% itemjoin 决定相邻条目之间插入的内容。
% itemjoin* 决定最后两个条目之间插入的内容，通常与 itemjoin 保持一致。
% 内容过长时，行内列表仍会自动换行。
\newlist{inlineenum}{enumerate*}{1}
\setlist[inlineenum,1]{
  label=(\alph*),
  labelsep=0.5em,
  itemjoin={；\quad},
  itemjoin*={；\quad}
}

% 其他常用 label 写法：
%   label=\arabic*.     -> 1. 2. 3.
%   label=(\arabic*)   -> (1) (2) (3)
%   label=\alph*)      -> a) b) c)
%   label=(\roman*)    -> (i) (ii) (iii)
%   label=\Alph*.      -> A. B. C.

% ============================================================
% 题目与回答环境
% ============================================================
% problem 使用独立计数器，每次进入环境时自动加一。
% \chinese{problem} 将题号显示为“一”“二”“三”等中文数字。
% \refstepcounter 使题目可以配合 \label 和 \ref 进行交叉引用。
% “题目一”等题目标识使用白底黑字描边框，与黑底白字的“回答”标签相区分。
\newcounter{problem}
\newenvironment{problem}
  {\refstepcounter{problem}%
   \par\medskip\noindent
   \begingroup
   \setlength{\fboxsep}{3pt}%
   \setlength{\fboxrule}{0.8pt}%
   \fbox{\bfseries 题目\chinese{problem}}%
   \endgroup
   \quad\ignorespaces}
  {\par\medskip}

% answer 环境用于书写答案。
% “回答”使用黑底白字标签，保证黑白打印时仍有醒目的视觉区分；
% 答案结束处会自动添加一条黑色细分隔线。
\newenvironment{answer}
  {\par\smallskip\noindent
   \begingroup
   \setlength{\fboxsep}{4pt}%
   \colorbox{black}{\textcolor{white}{\bfseries 回答}}%
   \endgroup
   \quad\ignorespaces}
  {\par\smallskip\noindent\rule{\linewidth}{0.4pt}\medskip}

\begin{document}

\begin{center}
  {\LARGE\bfseries \coursename\quad \assignmentnumber}\par
  \vspace{0.7em}
  \small
  教师：\teachername\qquad
  姓名：\studentname\qquad
  学号：\studentid
\end{center}
\vspace{0.4em}

% ============================================================
% 从这里开始写正文。
%
% 普通分行小问使用 enumerate：
%   \begin{enumerate}
%     \item ...
%     \item ...
%   \end{enumerate}
%
% 同行小问使用已全局配置好的 inlineenum：
%   \begin{inlineenum}
%     \item ...
%     \item ...
%   \end{inlineenum}
%
% 每道题的基本结构：
%   \begin{problem}
%     ...
%   \end{problem}
%   \begin{answer}
%     ...
%   \end{answer}
%
% 如需交叉引用，可在题目开头添加标签：
%   \begin{problem}\label{prob:example}
% 随后使用“第 \ref{prob:example} 题”引用该题。
% ============================================================

\begin{problem}
设函数 $f(x)=\mathrm{e}^{\mathrm{i}x}$。
\begin{inlineenum}
  \item 求 $f'(x)$
  \item 计算 $\displaystyle\int_0^{2\pi}f(x)\,\mathrm{d}x$
\end{inlineenum}。
\end{problem}

\begin{answer}
由指数函数的求导与积分公式，有
\begin{equation}
  f'(x)=\mathrm{i}\mathrm{e}^{\mathrm{i}x}.
\end{equation}
并且
\begin{equation}
  \int_0^{2\pi}f(x)\,\mathrm{d}x
  =\int_0^{2\pi}\mathrm{e}^{\mathrm{i}x}\,\mathrm{d}x
  =0.
\end{equation}
\end{answer}

\begin{problem}\label{prob:commutator}
已知正则对易关系 $[\hat{x},\hat{p}]=\mathrm{i}\hbar$，求 $[\hat{x},\hat{p}^{2}]$。
\end{problem}

\begin{answer}
利用恒等式 $[\hat{A},\hat{B}\hat{C}]=[\hat{A},\hat{B}]\hat{C}
+\hat{B}[\hat{A},\hat{C}]$，可得
\begin{equation}
  [\hat{x},\hat{p}^{\,2}]
  =[\hat{x},\hat{p}]\hat{p}+\hat{p}[\hat{x},\hat{p}]
  =2\mathrm{i}\hbar\hat{p}.
\end{equation}
\end{answer}

\end{document}
